\section*{Summary}
\addcontentsline{toc}{section}{Summary}

\begin{itemize}
  \item CAN does framing, length, checksum, collision handling and acknowledgement in silicon.
    What is left for the software is the identifier, the length, the data bytes --- and reading
    off how it went.
  \item CAN has \textbf{no destination address}. The identifier describes the content, and every
    node filters for itself.
  \item The controller signals with single-cycle pulses. \textbf{The register bank} turns them
    into sticky bits with an \textbf{explicit clear} --- and that clearing is the driver's
    responsibility.
  \item \code{TX_SEND} and \code{RX_ACK} require bit 0 to be set; \code{ERROR_FLAGS} is cleared
    only by an \emph{all-zero} word. Opposite rules, and a common bug.
  \item \code{TX_DATA_LO} carries data bytes \textbf{0--3}, with byte 0 in bits 31--24.
  \item Three limitations reach up into the code: a lost arbitration looks like a successful
    transmission, the error flag is one bit for four causes, and the receive side does not buffer.
  \item The hardware validates nothing. The validation lives in the driver, in exactly one place.
  \item Two seams: \code{Interface} makes the application testable, \code{ByteTransport} makes the
    driver testable.
\end{itemize}
